Des muons sont produits à \qty{\v{h}}{km} d'altitude par des rayons cosmiques et descendent verticalement vers le sol à $\num{\v{beta}}\,c$. Au repos, la durée de vie moyenne d'un muon est de \qty{2.2}{\micro\second}.

\begin{enumerate}[label=\alph*)]
    \item Selon un observateur au sol, combien de temps les muons mettent-ils pour atteindre le sol ?
    \item Combien de temps s'écoule selon une horloge qui voyagerait avec les muons ? Lequel de ces deux intervalles est un temps propre ?
    \item Le nombre de muons qui survivent après un temps propre $\tau$ est $N = N_0\,e^{-\tau/\tau_0}$ avec $\tau_0 = \qty{2.2}{\micro\second}$. Quelle fraction des muons atteint le sol ? Comparez avec la fraction qu'on obtiendrait sans relativité.
\end{enumerate}
